\documentclass[11pt]{article}

\usepackage[margin=1in]{geometry}
\usepackage{amsmath, amssymb}
\usepackage{listings}
\usepackage{xcolor}
\usepackage{hyperref}
\usepackage{enumitem}
\usepackage{titlesec}
\usepackage{parskip}
\usepackage{booktabs}

% Code listing style
\lstset{
    basicstyle=\ttfamily\small,
    keywordstyle=\color{blue}\bfseries,
    commentstyle=\color{gray}\itshape,
    stringstyle=\color{red},
    frame=single,
    breaklines=true,
    showstringspaces=false,
    tabsize=4
}

\title{\textbf{CS3210: Operating Systems} \\ Lecture 1 --- Introduction and Course Overview}
\author{John Kim \\ Georgia Institute of Technology}
\date{January 13, 2026}

\begin{document}

\maketitle
\tableofcontents
\newpage

% ============================================================
\section{What is an Operating System?}
% ============================================================

\subsection{Three Classical Responsibilities}

Before diving into the technical content of the course, it is worth asking a deceptively simple
question: \emph{what is an operating system?} The classical textbook answer is that an operating
system has three responsibilities: \textbf{managing the computer's resources}, \textbf{establishing
a user interface}, and \textbf{executing and providing services for applications}. While this
enumeration is broadly accurate, it is worth considering a richer perspective.

Professor Tumanov's preferred definition is more conceptually tight: \emph{an operating system
abstracts \textbf{inconvenient} low-level interfaces into \textbf{convenient} high-level
interfaces.} This framing places \emph{abstraction} at the center of OS design rather than
listing capabilities, and will remain the organizing lens for the entire course.

\subsection{Abstractions as the Core Mechanism}

Operating systems work by providing \textbf{abstractions} that hide the complexity and inconvenience
of the underlying hardware. Consider two canonical examples. First, a single physical CPU is
abstract\-ed into the illusion of many independent processes running concurrently---this is
\textbf{multiplexing}. Second, a single flat set of physical memory cells is made to look like
many independent, isolated address spaces, one per process. In both cases the OS interposes
between the hardware and the application, translating a difficult-to-use physical resource into
a convenient logical one.

A useful question to ask throughout the course is: what \emph{other} abstractions can you think
of? The lesson is that the pattern is universal and the design space is vast.

\subsection{Properties of OS Abstractions}

Good OS abstractions exhibit a cluster of desirable properties. The full list includes:

\begin{itemize}
    \item \textbf{Isolation} --- one process's faults or misbehavior cannot corrupt another's state.
    \item \textbf{Multiplexing} --- a single physical resource can be shared among multiple
          consumers simultaneously.
    \item \textbf{Naming} --- resources are given human-meaningful or process-meaningful names
          (e.g., file paths, process IDs) rather than raw hardware addresses.
    \item \textbf{Atomicity} --- certain operations are guaranteed to complete entirely or not at
          all, with no partial visible state in between.
    \item \textbf{Reliability / Fault Tolerance} --- the system continues to function correctly
          despite individual component failures.
    \item \textbf{Persistence} --- state can survive beyond the lifetime of any individual process
          or even a system reboot.
    \item \textbf{Resource Management} --- the OS enforces policies for sharing CPU time, memory,
          disk I/O, and network bandwidth fairly and efficiently across all running processes.
\end{itemize}

Not every abstraction exhibits all of these properties; part of OS design is deciding which
properties matter most for a given context and making the necessary trade-offs.

% ============================================================
\section{Managing Complexity Through Abstraction}
% ============================================================

A recurring theme in systems is that abstraction is the primary tool for managing complexity.
The lecture presented a three-layer model that captures how abstraction enables systems to scale:

\begin{enumerate}
    \item \textbf{Learning Algorithm / Common Patterns} (top layer) --- identifies recurring
          computation patterns in applications (e.g., map-reduce, iterative gradient descent,
          request--response).
    \item \textbf{Abstraction (API)} (middle layer) --- defines a \emph{narrow interface} through
          which the application expresses its intent without specifying implementation details.
    \item \textbf{System} (bottom layer) --- exploits the limited, well-defined interface to
          address system design challenges such as parallelism, data locality, networking,
          scheduling, fault tolerance, and stragglers.
\end{enumerate}

The power of the narrow API in the middle is that it allows the system below to change freely
(e.g., to improve performance, change scheduling policy, or tolerate faults) without requiring
changes to the applications above. This \emph{decoupling} is what makes large-scale systems
tractable.

\subsection{What Counts as an Operating System?}

This layered model raises an interesting boundary question: what counts as an ``operating
system''? The classic answer is the kernel. But Richard Stallman famously argued for a broader
view: the GNU project is the operating system, and Linux is merely one of its kernels. In his
formulation, the tools, libraries, and user-space infrastructure together constitute the OS.

For this course we will focus on the kernel layer---the ring-0 x86 software that directly manages
hardware. We use \textbf{xv6}, a small but complete teaching kernel developed at MIT, as our
primary vehicle for hands-on projects. xv6 is small enough to read and modify but rich enough to
illustrate all core OS mechanisms.

% ============================================================
\section{Course Overview}
% ============================================================

\subsection{Design as a First-Class Goal}

The course title is \emph{Design of Operating Systems}, and the emphasis on \emph{design} is
intentional. Many OS courses focus primarily on describing how existing systems work. This course
goes further: students are expected to understand \emph{why} particular abstractions exist,
\emph{what} trade-offs they encode, and how to \emph{design} systems from scratch given
conflicting constraints. As Prof.\ Tumanov put it, the goal is for students to be able to
``understand, express, and quantify the \textbf{trade-offs} in system design.''

\subsection{Course Structure}

The course is primarily \textbf{project-based}. Students work directly in the xv6 kernel---a real
x86 ring-0 environment---rather than a simulator. Approximately 55\% of the final grade is
determined by four lab projects; two midterm exams account for the remaining 40\%, with a small
supervised-lab component making up the last 3\%.

The lab projects are substantial. Prof.\ Tumanov estimates that each lab will take approximately
\textbf{30 hours}, with Lab 3 potentially requiring 40--50 hours. Students should plan for
roughly \textbf{10 hours per week} of lab work during active lab periods. The autograder accepts
only $k$ feedback-given submissions per day, so starting early is not merely advisable---it is a
structural requirement of the course. The autograder reports how \emph{complete} a solution is,
not \emph{what is wrong with it}; students must write their own unit tests to diagnose failures.

The grading breakdown is as follows:

\begin{center}
\begin{tabular}{ll}
\toprule
\textbf{Item} & \textbf{Weight} \\
\midrule
Lab 0 (XV6 orientation) & 2\% \\
Lab 1 (Startup, CPU stacks, debugging, physical memory) & 10\% \\
Lab 2 & 15\% \\
Lab 3 & 20\% \\
Lab 4 & 10\% \\
Supervised Labs & 3\% \\
Exam 1 & 20\% \\
Exam 2 & 20\% \\
\bottomrule
\end{tabular}
\end{center}

Note that the withdrawal deadline falls between Lab 2 and Lab 3, giving students a concrete
checkpoint at which to assess their standing before committing to the more demanding second half
of the course.

\subsection{Scope and Limitations}

Operating systems are enormous software artifacts---millions of lines of code, highly complex,
and far broader than just their kernels. No single course can cover the full breadth of useful OS
knowledge. The course explicitly covers the subset of OS concepts that (a) can be meaningfully
implemented in xv6, (b) illuminate important design trade-offs, and (c) will be tested on exams.
Students should come in with that expectation and use the course as a foundation for independent
exploration later.

% ============================================================
\section{Prerequisites and Course Policies}
% ============================================================

\subsection{Prerequisites}

Operating systems sit at a low level of the software stack and interact directly with hardware.
The course therefore requires a broad and solid background. The \emph{strict} prerequisites---
without which a student will struggle---are:

\begin{itemize}
    \item \textbf{C programming} --- all kernel code is written in C; pointer manipulation,
          memory allocation, and low-level bit operations are used throughout.
    \item \textbf{Linux shell} --- projects are developed and debugged on the command line;
          comfort with shell scripting, \texttt{make}, \texttt{gdb}, and related tools is assumed.
    \item \textbf{CS 2200 (Systems and Networks)} --- a thorough understanding of computer
          organization, instruction sets, system calls, and basic networking is required.
\end{itemize}

The \emph{recommended} (but not strictly required) prerequisites are CS 2110 (Computer
Organization and Programming) and CS 3220 (Processor Design). Students lacking these are still
welcome but should expect to work harder on the hardware-facing parts of the course.

A \textbf{diagnostic quiz} will be available on Gradescope during the first week (due Friday,
January 16, 2026). Its purpose is twofold: it lets the teaching staff calibrate supervised lab
topics to actual knowledge gaps, and it lets each student honestly assess their own preparedness
before the withdrawal deadline.

\subsection{Classroom Conduct and Participation}

The course is delivered \textbf{in-person only}. Lectures are not recorded as a default, though
select key lectures (e.g., Virtual Memory Management, Scheduling, File Systems, Atomicity) may
be recorded and shared through Canvas as exam preparation aids. Attendance is strongly encouraged
and considered essential even though it is not formally graded.

Students are encouraged to ask questions at any time during lecture. The only conduct norm is
mutual respect: questions should stem from genuine curiosity, not from a desire to signal prior
knowledge. All students, regardless of background, identity, or experience, are entitled to a
comfortable and inclusive learning environment.

\subsection{Academic Integrity}

Lab projects may be completed individually or in pairs (groups of 1--2). \emph{You may not share
code with anyone outside your partnership.} The distinction between permissible collaboration and
cheating is framed as follows: \textbf{cheating} is a one-sided exchange (one party benefits from
the other's work without reciprocal learning); \textbf{collaboration} is both parties discovering
and solving problems together. When in doubt, the correct action is to avoid sharing and ask the
course staff for clarification.

Georgia Tech's Honor Code applies in full. Students are expected to uphold the ideals of honor
and integrity in all work submitted for this course.

\subsection{Engineering Ethics and Impact}

The lecture closed with a broader observation about engineering responsibility. Competence in
systems design is not merely an academic exercise---systems that are poorly designed, or
credentials obtained dishonestly, can have real-world consequences that extend far beyond the
individual engineer. Students are expected to take both their learning and their integrity
seriously for this reason.

% ============================================================
\section{Upcoming Assignments}
% ============================================================

The following tasks are due in the first two weeks of the semester:

\begin{itemize}
    \item \textbf{Diagnostic Quiz} --- due Friday, January 16, 2026 (available on Gradescope).
          Also serves as a Phase II registration checkpoint.
    \item \textbf{Notecard} --- a short personal introduction for the instructor, due Friday
          January 16, 2026.
    \item \textbf{Lab 0} (Get to Know XV6) --- assigned January 16; due January 23, 2026.
    \item \textbf{Lab 1} (Startup, CPU Stacks, Debugging, Physical Memory) --- assigned January
          23; due February 6, 2026. Students have two weeks.
\end{itemize}

The first reading assignment is also due soon: James Mickens, ``The Night Watch,''
\url{https://www.usenix.org/system/files/1311_05-08_mickens.pdf}. Although ungraded, this essay
is required reading and sets a memorable tone for the kind of thinking this course values.

% ============================================================
\section{Summary}
% ============================================================

\begin{itemize}
    \item \textbf{Operating system definition}: An OS abstracts \emph{inconvenient} low-level
          hardware interfaces into \emph{convenient} high-level ones for applications.
    \item \textbf{Abstraction properties}: Key properties of OS abstractions include isolation,
          multiplexing, naming, atomicity, reliability, persistence, and resource management.
    \item \textbf{Three-layer model}: Systems manage complexity by identifying common patterns,
          defining a narrow API, and exploiting that API to address parallelism, scheduling,
          fault tolerance, and other system-level challenges.
    \item \textbf{Course design focus}: CS3210 emphasizes understanding and quantifying
          design trade-offs, not just describing how existing systems work.
    \item \textbf{xv6}: The primary project vehicle is xv6, a real x86 ring-0 kernel; all
          lab work involves modifying actual kernel code.
    \item \textbf{Grading}: 55\% labs (4 projects), 40\% exams (2), 3\% supervised labs;
          workload is approximately 30 hours per lab.
    \item \textbf{Integrity}: Collaboration between lab partners is encouraged; sharing code
          with non-partners is cheating. The Honor Code applies in full.
\end{itemize}

% ============================================================
\section{References}
% ============================================================

\begin{itemize}
    \item Tumanov, A. \textit{CS3210: Design of Operating Systems --- Lecture 1 Slides}.
          Georgia Institute of Technology. January 13, 2026.
    \item Mickens, J. ``The Night Watch.'' \textit{USENIX ;login: Logout}. November 2013.
          \url{https://www.usenix.org/system/files/1311_05-08_mickens.pdf}
    \item Silberschatz, A.; Galvin, P. B.; Gagne, G. \textit{Operating System Concepts}.
          10th ed. John Wiley \& Sons. 2018.
    \item Cox, R.; Kaashoek, F.; Morris, R. \textit{xv6: a simple, Unix-like teaching operating
          system}. MIT CSAIL. \url{https://pdos.csail.mit.edu/6.828/2022/xv6.html}
\end{itemize}

% ============================================================
\section*{Personal Notes}
% ============================================================
% Add your own notes below this line

\end{document}
